%======================================================================
\section{Open questions}\label{sec:open}
%======================================================================

\paragraph{The exact scale.} For the game of Theorem~A,
\[
c\,\eps^{-1}\bfun(\eps)\ \le\ \log\Ent(G,\eps)\ \le\ \log\Entme(G,\eps)\ \le\ C\eps^{-1}\log(1/\eps).
\]
The factor $\bfun(\eps)=2^{-O((\log\log(1/\eps))^2)}$ comes only from the pair areas $\Lambda_N$, that is from the Dehn function of $H_3$, and
the factor $\log(1/\eps)$ only from the number of affine hashes over which the models of the upper bound average. Under the polynomial bound~\ref{H-poly}, which
we do not use, $\bfun(\eps)$ improves to a power of $\log(1/\eps)$ (\cref{rem:sector-poly}).

\begin{question}\label{q:scale}
Is $\log\Ent(G,\eps)$ of order $\eps^{-1}\log(1/\eps)$, or of order $\eps^{-1}$? In $\GP$, is $\log\hlp(J;2,\eps)$ of order
$\eps^{-2}\log(1/\eps)$?
\end{question}

The lower bound sees $2^N\approx\eps^{-1}$ lamps on $N\approx\log(1/\eps)$ points, one qubit each. The upper bound keeps about $\eps^{-1}$
qubits as well, but averages over a family of about $2^{6mk}$ affine hashes, with $k\approx\log(1/\eps)$ bits, and the size of that family is
where the factor $\log(1/\eps)$ comes from. Deciding which count is right seems to need either a lower bound that also charges for the
points a lamp uses, or models built from a much smaller family of hashes.

\paragraph{The Dehn function of Houghton's group.} The proof of \cref{cor:dehn-qp} halves Lee's recursion for the commutator of two
transpositions at distance $M$ and obtains $M^{(1+o(1))\log M}$. The Dehn function of $H_3$ is at least quadratic: $H_3$ is not hyperbolic,
since its finitary permutations have unbounded finite orders, and a finitely presented group with subquadratic Dehn function is
hyperbolic~\cite{Bri02}.

\begin{question}\label{q:dehn}
Is the Dehn function of $H_3$ quadratic, improving the polynomial bound~\ref{H-poly} of our preprint~\cite{DM-H3}? In particular, do the commutators $[\sigma_0,\sigma_M]$ of far-apart transpositions have area $O(M^2)$?
\end{question}

A quadratic bound would give $\Lambda_N=O(N^4)$ and replace $\bfun(\eps)$ by $(\log(1/\eps))^{-O(1)}$; it would not touch the factor
$\log(1/\eps)$ of \cref{q:scale}.

\paragraph{Faster growth.} By $\Pi_2$-completeness~\cite{MNY22}, every total computable bound on $\Ent(G,\eps)$ in terms of the game and
the precision is exceeded by some game of quantum value one, but the argument exhibits no such game. The exponent $1$ of Theorem~A is the
natural limit of the present method: a strategy can only certify about $\eps^{-1}$ independent qubits
at precision $\sqrt\eps$ when the errors add in squares, so faster growth must come from a mechanism other than counting commuting Pauli
pairs, for instance from lamps whose state space itself grows with the precision.

\begin{question}\label{q:faster}
Is there an explicit two-player game $G$ of quantum value one such that $\log\log\Ent(G,\eps)\ge\eps^{-c}$ for some $c>0$? Can it be a
linear system game whose solution group is amenable?
\end{question}

For amenable solution groups the second part asks for a central involution $J$ whose state-weighted profile is at least doubly
exponential; by Theorem~D this is exactly what such a game needs.

\paragraph{Keeping the state.} By Theorem~D the entanglement of a linear system game is the state-weighted profile of $J$, and the
entanglement of maximally entangled strategies is the flat profile $\hlp(J;2,\cdot)$, at the same precision $\sqrt\eps$ up to constants.
In all examples we know, the logarithms of the two profiles agree up to factors $(1/e)^{o(1)}$ as the precision $e\to0$: for $\GP$ by Theorem~B, and for the group
$K$ of Slofstra and Vidick, where the profiles themselves agree up to constant factors, by Theorem~F.

\begin{question}\label{q:flat}
Is there a finitely presented group with a central involution $J$ for which $\log\hlps(J;e)$ is much smaller than
$\log\hlp(J;2,e)$, say $\log\hlps(J;e)\le(\log\hlp(J;2,e))^{1-c}$? If such a group can be the solution group of a linear system game of
quantum value one, then, by Theorem~D and~\cite[Cor.~4.6]{SV18}, strategies of that game far from maximally entangled save a power of the
entanglement.
\end{question}

\paragraph{A smaller game.} The game has $1687$ equations because it encodes, through Slofstra's embedding, an $E\Gamma_0$ presentation,
with $27$ generators, of a group that maps onto $\GPp$. Most of the equations are gadgets: each conjugacy relation of the presentation
costs fourteen equations. The constants of Theorem~A are explicit but astronomically small, and a smaller game would improve them too.

\begin{question}\label{q:small}
What is the least number of equations of a linear system game, or of questions of a two-player game, of quantum value one in which
$\Ent(G,\eps)$ grows faster than every polynomial in $1/\eps$?
\end{question}

\paragraph{Other measures of entanglement.} The lower bounds of Theorems~A and~E hold for the entanglement entropy of pure strategies and for
the entanglement of formation of mixed ones. For mixed states the entanglement of formation is at least the entanglement cost, the
distillable entanglement and the relative entropy of entanglement, so lower bounds on it do not bound these.

\begin{question}\label{q:distill}
Do near-optimal mixed strategies for CHSH$(n)$, or for the game of Theorem~A, have distillable entanglement of the same order as their
entanglement of formation?
\end{question}
